\documentclass[10pt,a4paper]{amsart}
\usepackage[margin=19mm]{geometry}
\usepackage{amsmath,amssymb,mathtools}
\usepackage[scr=rsfs,cal=euler]{mathalfa}
\usepackage{xfrac}
\usepackage[unicode,hidelinks,bookmarks=false]{hyperref}
\newcommand{\ie}{i.e.\ }
\newcommand{\cf}{cf.\ }
\newcommand{\resp}{respectively\ }
\newcommand{\Subsection}{\subsection}
\providecommand{\nobreakdash}{\nobreak\hbox{-}}
\setlength{\emergencystretch}{2em}
\multlinegap0pt
\usepackage{bm}
\usepackage{bbm}
\usepackage{dsfont}
\usepackage{xspace}
\usepackage{cancel}
\usepackage{mathtools}
\usepackage[shortlabels]{enumitem}
\usepackage{amsthm}
\makeatletter
\@ifundefined{theorem}{\newtheorem{theorem}{Theorem}}{}
\@ifundefined{lemma}{\newtheorem{lemma}[theorem]{Lemma}}{}
\@ifundefined{proposition}{\newtheorem{proposition}[theorem]{Proposition}}{}
\makeatother
\DeclareMathOperator{\CB}{\mathcal{CB}}
\DeclareMathOperator{\amconv}{amconv}
\newcommand{\N}{{\mathbb{N}}}
\newcommand{\n}[1]{ \left\|#1\right\| }
\title[Revisiting Operator $p$-Compact Mappings]{Revisiting Operator $p$-Compact Mappings}
\author{Javier Alejandro Ch\'avez-Dom\'inguez, Ver\'onica Dimant, Daniel Galicer}
\date{Source: arXiv:2405.18571v2 (CC BY 4.0)}
\begin{document}
\maketitle

\noindent\emph{Source and licence.} Revisiting Operator $p$-Compact Mappings. Version arXiv:2405.18571v2 (CC BY 4.0), distributed under the Creative Commons Attribution 4.0 International licence, as recorded in the arXiv licence metadata for this exact version and shown on the abstract page https://arxiv.org/abs/2405.18571v2. Full rights notice: licenses/arxiv-2405.18571-CC-BY-4.0.txt.\par\smallskip

\noindent\textit{Editorial scope.} Counted unit: the Theorem whose source label is \texttt{thm:equivalences-p-compact-u} in the article (source0.tex lines 1632--1640, source label \texttt{thm:equivalences-p-compact-u}), delivered together with the article's own complete original proof of it (source0.tex lines 1642--1668, one \texttt{proof} environment of 27 lines). No mathematics is written, repaired or re-derived: statement and proof are the article's own text line for line. Required context retained uncounted: lemma-p-compact-closure (lines 359--362); lemma-amconv-preserves-p-compactness (lines 1540--1542); lemma:abs-matrix-convex-properties (lines 1555--1562); prop-matrix-unit-ball-of-ell1-sum (lines 1607--1610). Every definition, display and label of those lines is retained unchanged. Omitted from this package and not redistributed: 1--358, 363--1539, 1543--1554, 1563--1606, 1611--1631, 1641--1641, 1669--1981. Nothing inside a retained excerpt is omitted or altered: every retained body is delivered exactly as the article's lines are written, so no omission receipt of any kind accompanies this package. Citations: the retained excerpts carry no citation command at all, so no bibliography entry of the article is transcribed and no reference is redistributed. Reference adaptation: none was needed and none was made; every reference inside the delivered unit resolves inside it, so no unresolved reference or citation is produced. Printed slips kept literal: no printed word, symbol, hypothesis or number of the counted statement or of its proof was altered. Layout and class: the article's own class is \texttt{amsart} with packages including amssymb, amsmath, amsthm, verbatim, enumerate, graphicx, epsfig, color, soul, xy, wasysym, url, mathrsfs, cancel; the wrapper is the lane's pinned \texttt{amsart} editorial wrapper at 10pt on A4, which restates the theorem-like environments the retained text needs and adds the article's own macro definitions that the retained text needs (\texttt{\textbackslash CB}, \texttt{\textbackslash N}, \texttt{\textbackslash amconv}, \texttt{\textbackslash n}), with extra wrapper packages (bm, bbm, dsfont, xspace, cancel, mathtools, enumitem). The delivered numbering is the unit's own. Assets: no retained span contains a figure environment, a graphics-inclusion command, a PDF-page inclusion, a table or a TikZ picture; no image file is copied into this package, the package declares no asset dependency and no image file is redistributed with it. Licence: version arXiv:2405.18571v2 of this article is distributed under the Creative Commons Attribution 4.0 International licence, as recorded in the arXiv licence metadata for this exact version and shown on the abstract page https://arxiv.org/abs/2405.18571v2. A full rights notice is delivered at licenses/arxiv-2405.18571-CC-BY-4.0.txt. Characters: every retained excerpt is pure ASCII, so the delivered engine, XeLaTeX with its default Latin Modern text font, reports no missing character.
\begin{lemma}\label{lemma-p-compact-closure}
Let $1 \le p \le \infty$.
A matrix set $\mathbf{K}$ over an operator space $V$ is relatively operator $p$-compact if and only if so is $\overline{\mathbf{K}}$. Furthermore, $\frak m_p^o(\mathbf{K}) = \frak m_p^o(\overline{\mathbf{K}})$.
\end{lemma}

\begin{lemma}\label{lemma-amconv-preserves-p-compactness}
Let $\mathbf{K}=(K_n)_n$ be a matrix set over $V$, and $1\le p \le \infty$. Then $\mathbf{K}$ is relatively operator $p$-compact if and only if so is $\amconv(\mathbf{K})$. Moreover, in this case $\frak m_p^o(\mathbf{K}) =\frak m_p^o(\amconv(\mathbf{K}))$.
\end{lemma}

\begin{lemma}\label{lemma:abs-matrix-convex-properties}
Let $V$ be an operator space and $\mathbf{K} = (K_n)$ an absolutely matrix convex  set over $V$.
\begin{enumerate}[(a)]
\item If $\rho : M_n \to M_m$ is a complete contraction, then $(\rho \otimes Id_V) K_n \subseteq K_m$.
\item If $\xi \in M_m(\mathcal{S}_1^n)$ with $\n{\xi} \le 1$ and $x \in K_n$, then
$(\Theta^x)_m(\xi) \in K_m$.
\end{enumerate}
\end{lemma}

\begin{proposition}\label{prop-matrix-unit-ball-of-ell1-sum}
Let $\{V_i\}_{i \in I}$ be a collection of operator spaces. For each $i_0\in I$ let $J_{i_0} : V_{i_0} \to \ell_1( \{ V_i\}_{i \in I} )$ be the canonical complete isometry, that is, the one that sends $v \in V_{i_0}$ to the vector with $v$ in the $i_0$-th position and 0 everywhere else.
Then the closed matrix unit ball of $\ell_1( \{ V_i\}_{i \in I} )$ is the closure of the absolutely matrix convex hull of $\bigcup_{i\in I} J_i(\mathbf{B}_{V_i})$.
\end{proposition}

\begin{theorem}\label{thm:equivalences-p-compact-u}
Let $\mathbf{K}$ be a completely bounded matrix set over $V$.
The following are equivalent:
\begin{enumerate}[(i)]
\item $\mathbf{K}$ is relatively operator $p$-compact.
\item $u_\mathbf{K} : \ell_1\big( \{ \mathcal{S}^{n_x}_1 \}_{x \in \mathbf{K}}\big) \to V$ is operator $p$-compact.
\end{enumerate}
Moreover, in this case $\frak m_p^o(\mathbf{K}) = \kappa_p^o(u_\mathbf{K})$.
\end{theorem}

\begin{proof}
For $y\in \mathbf{K}$, denote by $J_y : \mathcal{S}_1^{n_y} \to \ell_1( \{ \mathcal{S}^{n_x}_1 \}_{x \in \mathbf{K}}) $ the canonical inclusion. It is obvious from the definition of $u_{\mathbf{K}}$ that for each such $y$ we have $u_{\mathbf{K}}J_y = \Theta^y$.

$(ii) \Rightarrow (i)$:
Fix $y\in \mathbf{K}$.
Let $\xi \in M_{n_y}(\mathcal{S}_1^{n_y}) \equiv \CB(M_{n_y},M_{n_y})$ be the matrix of norm one that corresponds to the identity map $M_{n_y} \to M_{n_y}$, which is the matrix $\xi = (E_{kl})_{k,l=1}^{n_y}$ of matrix units.
Since $\Theta^yE_{kl} = y_{kl}$ for all $1 \le k,l \le n_y$, we have that $y = (u_{\mathbf{K}}J_y)_{n_y}(\xi)$.
Since $J_y$ is a complete isometry we know $J_y(\mathbf{B}_{\mathcal{S}_1^{n_y}}) \subseteq \mathbf{B}_{\ell_1( \{ \mathcal{S}^{n_x}_1 \}_{x \in \mathbf{K}}) }$, and therefore we have proved
$$
\mathbf{K} \subseteq u_{\mathbf{K}}\big( \mathbf{B}_{\ell_1( \{ \mathcal{S}^{n_x}_1 \}_{x \in \mathbf{K}}) } \big).
$$

Now, using that $u_{\mathbf{K}}$ is operator $p$-compact we conclude $\mathbf{K}$ is relatively operator $p$-compact and $\frak m_p^o(\mathbf{K}) \le \kappa_p^o(u_\mathbf{K})$.

$(i) \Rightarrow (ii)$:
For $x\in \mathbf{K}$, $m\in\N$ and $\xi \in B_{M_m(\mathcal{S}_1^{n_x})}$, note that $(u_{\mathbf{K}}J_x)_m(\xi) = (\Theta^x)_m(\xi)$
which belongs to $\amconv(\mathbf{K})$ by Lemma \ref{lemma:abs-matrix-convex-properties}.
This shows that $u_{\mathbf{K}} \Big( \bigcup_{x\in \mathbf{K}} J_x(\mathbf{B}_{\mathcal{S}_1^{n_x}}) \Big) \subseteq \amconv(\mathbf{K})$, from where it follows by linearity of $u_{\mathbf{K}}$ that
$u_{\mathbf{K}} \big( \amconv \Big( \bigcup_{x\in \mathbf{K}} J_x(\mathbf{B}_{\mathcal{S}_1^{n_x}}) \Big) \Big)\subseteq \amconv(\mathbf{K})$ and therefore, since $u_\mathbf{K}$ is continuous,
\[
u_{\mathbf{K}} \Big( \overline{\amconv \Big( \bigcup_{x\in \mathbf{K}} J_x(\mathbf{B}_{\mathcal{S}_1^{n_x}})} \Big) \Big) \subseteq
\overline{u_{\mathbf{K}} \Big( \amconv \Big( \bigcup_{x\in \mathbf{K}} J_x(\mathbf{B}_{\mathcal{S}_1^{n_x}}) \Big) \Big)}\subseteq \overline{\amconv(\mathbf{K})}.
\]
The desired conclusion now follows from Lemmas \ref{lemma-p-compact-closure} and 
 \ref{lemma-amconv-preserves-p-compactness},
and Proposition \ref{prop-matrix-unit-ball-of-ell1-sum}.
\end{proof}

\end{document}
